\scanpage{67}
\[
\unsure{v|\nu}: C\longrightarrow X\times Y,
\]
où $C$ est une courbe lisse projective, donc par
\[
v_1:C\unsure{\mathbin{-}|\to}X,\qquad v_2:C\mathbin{-}Y
\]
Il suffit alors de prendre le composé
\diagram{\begin{tikzcd}[ampersand replacement=\&]
B_X\arrow[r,dash,"v_1^*"] \& B_C\arrow[r,"\text{dualité}"] \& A_C\arrow[r,"v_{2*}"] \& A_Y
\end{tikzcd}}
ou encore, en termes de classes de correspondances divisorielles sur $B_X\times B_Y$, il suffit de composer la classe \deleted{\unsure{du diviseur|dualisante}} canonique sur $B_C\times B_C$, avec les hom $B_X\xrightarrow{v_1^*}B_C$, $B_Y\overset{v_2^*}{\mathbin{-}}B_C$. Ceci posé, si $X=Y$, $v_1=v_2$, on voit que \deleted{cette classe de} les conditions suivantes sont équivalentes :

(i) L'hom $B_X\mathbin{-}A_X=B_X^*$ est un hom de polarisation

(ii) L'hom $B_X\underset{\unsure{\cdot|\prime}}{\mathbin{-}}B_X^*=A_X$ est une isogénie

(iii) L'hom $\unsure{v|w}^*:B_X\mathbin{-}B_C$ est injectif \inserted{mod groupes finis.}

(iv) L'hom $v_*:A_C\mathbin{-}A_X$ est surjectif
% archive sheet number: 66

